\section{Forecasting Design}\label{sec:design}

\subsection{Models}

All models are linear regressions of the target \eqref{eq:target} on
predictors dated $t$, estimated stock by stock. The baseline, model $A$, is
the HAR model of \citet{corsi2009simple} in logs with a return term,
\begin{equation}
y_{i,t,h} = \beta_0 + \beta_d \ln RV^{(1)}_{i,t} + \beta_w \ln RV^{(5)}_{i,t}
+ \beta_m \ln RV^{(22)}_{i,t} + \gamma_1 r_{i,t} + \gamma_2 \lvert r_{i,t}\rvert
+ \varepsilon_{i,t,h},
\label{eq:har}
\end{equation}
where $RV^{(k)}_{i,t}$ is the mean Parkinson variance over days $t-k+1$ to
$t$, and $r_{i,t}$ is the return on day $t$. HAR-X, model $A_1$, adds the
levels of the VIX and MOVE indices at the close of day $t$. The persistence
blocks are added to HAR one at a time and then jointly. Three further
specifications, fixed before their results were computed, add a single block
to HAR-X, so that the comparison with HAR-X isolates that block.
Table~\ref{tab:models} lists the predictors.

\begin{table}[H]
\caption{Model specifications.}
\label{tab:models}
\footnotesize
\begin{tabularx}{\textwidth}{lX}
\toprule
Model & Predictors \\
\midrule
$A$ (HAR) & $\ln RV^{(1)}, \ln RV^{(5)}, \ln RV^{(22)}, r_t, \lvert r_t\rvert$ \\
$A_1$ (HAR-X) & HAR + VIX, MOVE \\
$A_2$ & HAR + own-stock $\hat d_{i,t}$, its weekly change, its standard deviation and slope over four weeks, own Hurst exponent and its weekly change \\
$A_3$ & HAR + cross-sectional mean $\bar d_t$ and dispersion $\sigma_{d,t}$ \\
$A_4$ & HAR + leave-one-out sector mean of $\hat d_{i,t}$ \\
$A_5$ & HAR + $\hat d_{i,t}$, VIX, MOVE, $\hat d_{i,t}\times$VIX, $\hat d_{i,t}\times$MOVE \\
$C$ & Union of $A_1$ to $A_5$ (18 predictors) \\
\midrule
\multicolumn{2}{l}{\textit{Registered HAR-X extensions}} \\
HAR-X + cross-sectional state & HAR-X + $\bar d_t$, $\sigma_{d,t}$ \\
HAR-X + sector state & HAR-X + leave-one-out sector mean of $\hat d_{i,t}$ \\
HAR-X + state $\times$ HAR terms & HAR-X + $\bar d_t$ and $\bar d_t\times \ln RV^{(k)}_{i,t}$, $k=1,5,22$ \\
\bottomrule
\end{tabularx}
\noindent{\footnotesize{\textit{Notes:} All predictors are dated at the close
of the origin day $t$. Rolling estimates use the 750 trading days ending on
day $t$. The last specification lets the persistence state change the
weights on the HAR components, the channel through which slower decay of
shocks would act.}}
\end{table}

\subsection{Estimation}

Each model is re-estimated by ordinary least squares at every weekly origin
on an expanding window of that stock's past observations. Because targets
overlap for $h>1$, a past observation enters the regression only if its
target window has closed by the origin: the observation for an earlier
origin $s$ is used at origin $t$ only if its target window, days $s+1$ to
$s+h$, ends no later than day $t$. Without this restriction the last
observations of each training sample contain variance that has not yet been
realized; at $h=22$ this inflated the apparent gain of the richer models in
an earlier version of this study. The first 431 origins (up to 14 June 2013)
form the initial training sample, and all models are evaluated over the same
646 origins, from 17 June 2013 to 14 April 2026, on the same (date, stock)
cells.

\subsection{Evaluation}

The loss function is the squared error of the log-variance forecast, pooled
over stocks and dates; Appendix~\ref{app:extra} reports the QLIKE loss for
variance forecasts corrected point-in-time for the log transformation.
Comparisons are based on the cross-sectional mean loss differential per date,
$\bar\delta_t = N^{-1}\sum_i (L^{(1)}_{i,t} - L^{(2)}_{i,t})$, which absorbs
the strong cross-sectional correlation of forecast errors. We report three
tests on $\bar\delta_t$, each with a Newey--West variance and bandwidth
$\max(\lceil h/5\rceil - 1, \lfloor 4(T/100)^{2/9}\rfloor)$, which is six
weeks for $T=646$:

\begin{itemize}
\item the Diebold--Mariano test of equal mean loss
\citep{diebold1995comparing} with the small-sample correction of
\citet{harvey1997testing};
\item for linear models that nest HAR-X, the \citet{clark2007approximately}
test, which adjusts the loss differential for the noise introduced by
estimating parameters whose population value is zero, and so tests whether
the extra predictors have nonzero population coefficients;
\item the conditional predictive ability test of \citet{giacomini2006tests},
with the constant and the most recent observed differential
$\bar\delta_{t-\lceil h/5\rceil}$ as instruments, which asks whether the
estimated model, with its estimation error, forecasts better.
\end{itemize}

\subsection{Registration}

The study was conducted in stages. Before the results of each stage were
computed, its specifications, comparisons, tests and decision rules were
recorded in a dated experiment ledger. Under the rule for the registered
HAR-X extensions, a specification adds information beyond HAR-X at a given
horizon only if its Holm-adjusted \citep{holm1979simple} Clark--West
$p$-value is below 0.05 and its mean squared error is lower than that of
HAR-X. After the results were known, we added the requirement that the gain
be at least 0.5\%; we report it separately because it was not registered. A
complete audit of the pipeline preceded the final run. Every correction was
registered before the corrected run and is listed in the ledger, which is
published with the code (Appendix~\ref{app:ledger}).
